\subsection{Prompt Construction}
\label{sec:prompt}

\noindent
Every LLM call in \toolname\ follows a fixed template, but the concrete prompt is specialized to the input program. A lightweight static classifier inspects the source code and assigns each program to one of five routes: recursive data structure, recursive program, array, numeric, and plain, where plain is the default route for programs that do not need helper definitions. The selected route determines which worked examples, category-specific ACSL guidance, and route-specific rule fragments are inserted on top of the universal ACSL pitfalls block. Thus, the same template yields substantively different prompts for, say, a numeric recurrence and a linked-list cursor.

This routing drives the five LLM call sites in the pipeline: loop invariant generation, precondition generation, postcondition/assigns generation, syntax repair, and verifier-guided semantic refinement. In addition, the symbolic executor supplies the state-dependent content of each prompt --- the variable state and the path constraints --- together with an annotation skeleton to be completed by the LLM. At refinement time, the verifier's typed error categories are additionally folded into the prompt. All prompt templates are listed in the supplementary material.
